\documentclass[UTF8,zihao=-4]{ctexart}
\usepackage[a4paper,margin=2.1cm]{geometry}
\usepackage{hyperref}
\usepackage{array}
\usepackage{longtable}
\usepackage{verbatim}
\hypersetup{hidelinks}
\setlength{\parindent}{0pt}
\setlength{\parskip}{0.55em}
\emergencystretch=4em
\sloppy
\title{05-程序分析与符号执行复习讲义}
\author{}
\date{}
\begin{document}
\maketitle
\setcounter{tocdepth}{1}
\tableofcontents
\newpage
\section{05 程序分析与符号执行复习讲义}


对应资料：


\begin{itemize}
\item \texttt{Slides/Slides/Ch9-1 Dataflow Analysis.pdf}
\item \texttt{Slides/Slides/Ch9-2-Symbolic Execution.pdf}
\item \texttt{Slides/Slides/Ch12-1-Interprocedure Analysis.pdf}
\item \texttt{Slides/Slides/Ch12-2 Pointer Analysis.pdf}
\item \texttt{Slides/Slides/handout/5 Dataflow Analysis.pdf}
\item \texttt{Slides/Slides/handout/6 Symbolic Execution.pdf}
\item \texttt{Slides/Slides/handout/7 Interprocedure Analysis.pdf}
\item \texttt{Slides/Slides/handout/8 Pointer Analysis.pdf}
\item \texttt{Review/compile.pdf}
\item \texttt{Exam/2025.pdf} 符号执行与指针别名题
\end{itemize}


\subsection{一、考试目标}


本章真题重点集中在：


\begin{quote}
\footnotesize
\begin{verbatim}
数据流分析：方向、meet、transfer
敏感性定义：path / flow / context / field
符号执行：符号状态、路径条件、约束求解
SAT：CNF、Tseitin、equisatisfiable
指针分析：Andersen 约束与闭包
\end{verbatim}
\end{quote}


零基础顺序：


\begin{quote}
\footnotesize
\begin{verbatim}
先背定义，再会写公式，最后做 Andersen 闭包。
\end{verbatim}
\end{quote}


\subsection{二、数据流分析框架}


数据流分析在 CFG 上反复传播事实。


基本记号：


\begin{quote}
\footnotesize
\begin{verbatim}
in[B]   : 基本块 B 入口处的事实
out[B]  : 基本块 B 出口处的事实
meet    : 多个前驱或后继的合并方式
transfer: 基本块内部对事实的改变
\end{verbatim}
\end{quote}


前向分析：


\begin{quote}
\footnotesize
\begin{verbatim}
in[B]  = meet(out[P]) for P in pred(B)
out[B] = transfer_B(in[B])
\end{verbatim}
\end{quote}


后向分析：


\begin{quote}
\footnotesize
\begin{verbatim}
out[B] = meet(in[S]) for S in succ(B)
in[B]  = transfer_B(out[B])
\end{verbatim}
\end{quote}


Worklist 模板：


\begin{quote}
\footnotesize
\begin{verbatim}
初始化 in/out
把相关基本块放入 worklist
while worklist 非空:
  取出 B
  根据方向重新计算 B 的 in/out
  若结果变化，把受影响的后继或前驱加入 worklist
\end{verbatim}
\end{quote}


\subsection{三、三种经典数据流分析}


可达定值 Reaching Definitions：


\begin{quote}
\footnotesize
\begin{verbatim}
问题：哪些定义能到达当前位置。
方向：前向
meet：union
公式：out[B] = gen[B] union (in[B] - kill[B])
初值：emptyset
\end{verbatim}
\end{quote}


可用表达式 Available Expressions：


\begin{quote}
\footnotesize
\begin{verbatim}
问题：哪些表达式在所有路径上已经计算过，且操作数未被改写。
方向：前向
meet：intersection
公式：out[B] = gen[B] union (in[B] - kill[B])
初值：全集，入口 out 为空集
\end{verbatim}
\end{quote}


活跃变量 Live Variables：


\begin{quote}
\footnotesize
\begin{verbatim}
问题：哪些变量当前值在之后还会被使用。
方向：后向
meet：union
公式：in[B] = use[B] union (out[B] - def[B])
初值：emptyset
\end{verbatim}
\end{quote}


\subsection{四、敏感性定义}


直接背：


\begin{quote}
\footnotesize
\begin{verbatim}
flow-sensitive:
区分程序点和语句顺序。

path-sensitive:
区分不同控制流路径。

context-sensitive:
区分函数调用上下文。

field-sensitive:
区分对象的不同字段。
\end{verbatim}
\end{quote}


答题补一句：


\begin{quote}
\footnotesize
\begin{verbatim}
敏感性越强，分析精度越高，状态数量和计算成本越高。
\end{verbatim}
\end{quote}


\subsection{五、符号执行}


符号执行把输入当成符号，而不是具体值。


过程：


\begin{quote}
\footnotesize
\begin{verbatim}
1. 输入变量赋符号值，例如 x = alpha_x。
2. 赋值语句更新符号状态。
3. 分支语句拆分路径。
4. 每条路径维护 path condition。
5. 用约束求解器判断 path condition 是否可满足。
\end{verbatim}
\end{quote}


例：


\begin{quote}
\footnotesize
\begin{verbatim}
if x > 0:
  y = x + 1
else:
  y = -x
\end{verbatim}
\end{quote}


路径 1：


\begin{quote}
\footnotesize
\begin{verbatim}
PC = alpha_x > 0
y = alpha_x + 1
\end{verbatim}
\end{quote}


路径 2：


\begin{quote}
\footnotesize
\begin{verbatim}
PC = alpha_x <= 0
y = -alpha_x
\end{verbatim}
\end{quote}


优点：


\begin{itemize}
\item 路径敏感。
\item 能生成测试输入。
\item 能发现深层路径错误。
\end{itemize}


成本：


\begin{itemize}
\item 路径爆炸。
\item 循环和递归需要限制或归纳。
\item 外部函数需要建模。
\end{itemize}


\subsection{六、CNF 与 SAT}


文字：


\begin{quote}
\footnotesize
\begin{verbatim}
x 或 not x
\end{verbatim}
\end{quote}


子句：


\begin{quote}
\footnotesize
\begin{verbatim}
(x or not y or z)
\end{verbatim}
\end{quote}


CNF：


\begin{quote}
\footnotesize
\begin{verbatim}
子句的合取。
(x or y) and (not x or z)
\end{verbatim}
\end{quote}


可满足：


\begin{quote}
\footnotesize
\begin{verbatim}
存在一组布尔赋值，使公式为真。
\end{verbatim}
\end{quote}


同时可满足 equisatisfiable：


\begin{quote}
\footnotesize
\begin{verbatim}
两个公式要么都可满足，要么都不可满足。
它不要求两个公式在同一变量集上逻辑等价。
\end{verbatim}
\end{quote}


\subsection{七、Tseitin 变换}


目的：


\begin{quote}
\footnotesize
\begin{verbatim}
把任意布尔公式线性规模转成 CNF，并保持同时可满足。
\end{verbatim}
\end{quote}


步骤：


\begin{quote}
\footnotesize
\begin{verbatim}
1. 给每个子表达式引入新变量。
2. 为新变量与子表达式的语义关系写 CNF。
3. 加入根变量为真的子句。
\end{verbatim}
\end{quote}


常用转换：


\begin{quote}
\footnotesize
\begin{verbatim}
x <-> (a and b)
(not x or a) and (not x or b) and (x or not a or not b)

x <-> (a or b)
(x or not a) and (x or not b) and (not x or a or b)

x <-> not a
(not x or not a) and (x or a)
\end{verbatim}
\end{quote}


真题证明模板：


\begin{quote}
\footnotesize
\begin{verbatim}
F1 = (a0 or ... or an or c) and (b0 or ... or bm or not c)
F2 = (a0 or ... or an or b0 or ... or bm)

若 F1 可满足：
  若 c 为 false，则第一项要求某个 ai 为 true。
  若 c 为 true，则第二项要求某个 bj 为 true。
  因此 F2 可满足。

若 F2 可满足：
  若某个 ai 为 true，令 c = true，则 F1 两个子句都为真。
  若某个 bj 为 true，令 c = false，则 F1 两个子句都为真。
  因此 F1 可满足。

所以 F1 与 F2 同时可满足。
\end{verbatim}
\end{quote}


\subsection{八、DPLL 与 CDCL}


DPLL：


\begin{quote}
\footnotesize
\begin{verbatim}
选择变量赋值
单子句传播
发现冲突则回溯
直到找到可满足赋值或证明不可满足
\end{verbatim}
\end{quote}


CDCL：


\begin{quote}
\footnotesize
\begin{verbatim}
在 DPLL 基础上加入冲突分析和学习子句。
学习子句用于剪枝，避免重复进入同类冲突。
\end{verbatim}
\end{quote}


\subsection{九、Andersen 指针分析}


Andersen 是包含约束的指针分析。


points-to 集合：


\begin{quote}
\footnotesize
\begin{verbatim}
pts(p) = p 能指向的对象集合
\end{verbatim}
\end{quote}


四类约束必须会写：


\begin{quote}
\footnotesize
\begin{verbatim}
y = &x:
  x in pts(y)

y = x:
  pts(x) subset pts(y)

*y = x:
  for v in pts(y), pts(x) subset pts(v)

y = *x:
  for v in pts(x), pts(v) subset pts(y)
\end{verbatim}
\end{quote}


\subsection{十、Andersen 真题最小例子}


题面结构：


\begin{quote}
\footnotesize
\begin{verbatim}
p = &a;
q = &b;
*p = q;
r = &c;
s = p;
t = *p;
*s = r;
\end{verbatim}
\end{quote}


逐句约束：


\begin{quote}
\footnotesize
\begin{verbatim}
a in pts(p)
b in pts(q)
for v in pts(p), pts(q) subset pts(v)
c in pts(r)
pts(p) subset pts(s)
for v in pts(p), pts(v) subset pts(t)
for v in pts(s), pts(r) subset pts(v)
\end{verbatim}
\end{quote}


闭包推导：


\begin{quote}
\footnotesize
\begin{verbatim}
pts(p) = {a}
pts(q) = {b}
由 *p = q 得 pts(a) 包含 {b}
pts(r) = {c}
由 s = p 得 pts(s) 包含 {a}
由 t = *p 且 pts(p) = {a} 得 pts(a) subset pts(t)，所以 pts(t) 包含 {b}
由 *s = r 且 pts(s) = {a} 得 pts(a) 包含 {c}
pts(a) = {b,c}
pts(t) = {b,c}
\end{verbatim}
\end{quote}


\subsection{十一、自测题}


题 1：活跃变量分析是前向还是后向？


答案：


\begin{quote}
\footnotesize
\begin{verbatim}
后向。
\end{verbatim}
\end{quote}


题 2：可用表达式的 meet 是什么？


答案：


\begin{quote}
\footnotesize
\begin{verbatim}
intersection。
\end{verbatim}
\end{quote}


题 3：符号执行为什么路径敏感？


答案：


\begin{quote}
\footnotesize
\begin{verbatim}
它为每条控制流路径维护独立的 path condition。
\end{verbatim}
\end{quote}


题 4：\texttt{y = *x} 的 Andersen 约束怎么写？


答案：


\begin{quote}
\footnotesize
\begin{verbatim}
for v in pts(x), pts(v) subset pts(y)
\end{verbatim}
\end{quote}

\end{document}
